\section{Existing Systems}
Conventional professor feedback systems in higher education are commonly form-based and batch-processed at semester end. While these systems collect basic ratings, they usually do not provide interactive dashboards, faculty engagement tracking, or integrated communication workflows.

\section{Identified Gaps}
\begin{enumerate}[label=\Alph*., leftmargin=1.5em]
    \item delayed analysis due to offline or end-semester processing,
    \item limited administrative visibility into ongoing faculty performance,
    \item low integration between rating systems and student support workflows,
    \item no mechanism for real-time trend visualization and ranking,
    \item no student-centric academic prediction support in the same platform.
\end{enumerate}

\section{Proposed Contribution}
EduRate addresses these gaps through a unified, role-based web platform that integrates:
\begin{enumerate}[label=\Alph*., leftmargin=1.5em]
    \item secure login and authorization for multiple stakeholders,
    \item professor ratings with actionable textual feedback,
    \item doubt posting and response workflow with usage control,
    \item performance analytics for both student and administrative views,
    \item email-based faculty reporting for institutional communication.
\end{enumerate}

\section{Comparative Perspective}
Compared to standard feedback portals, EduRate introduces continuous analytics, engagement features, and administrative action support. This improves the practical utility of feedback data and transforms it from static records into a decision-support resource.

\section*{Chapter Summary}
This chapter reviewed existing feedback system patterns, identified major operational and analytical gaps, and justified the contribution of EduRate as an integrated, role-aware, and decision-support-oriented platform.
